\scanpage{85}
% archive foliation: 87
revient à se donner un scindage de l'extension
\diagram{\begin{tikzcd}[ampersand replacement=\&,column sep=small]
(\ast)\qquad 1 \arrow[r] \& V \arrow[r] \& \Gamma \arrow[r,dash] \arrow[d,equal] \& \operatorname{Aut}(V) \arrow[r,dash] \arrow[d,"\simeq",draw=none] \& 1 \\
\&\& \operatorname{Nor}_G(V) \& \unsure{\mathfrak S}_J \&
\end{tikzcd}}
En termes géométriques, ces scindages sont associés
aux quatre systèmes de \inserted{paires de} faces antipodiques sur l'
octaèdre. \deleted{Ainsi, les} L'octaèdre combinatoire (non orienté
a priori) se définit à partir de l'ens $S$ de
ses sommets, et l'antipodisme \deleted{sur} $S$, de la façon
habituelle; ses faces sont les parties à trois éléments \deleted{qui}
$S'$ telles que $S'\cup\uline{a}(S')=S$. Pour un tel triplet,
le sous-groupe de $\Gamma$ qui l'invarie est un scindage
de l'extension $(\ast)$.

Le groupe $\Gamma$ est le groupe des autom. de
$S$ qui conserve \deleted{\unsure{des}} structure octaédrique
\uline{orientée} (définie par l'ens \inserted{$F\subset\unsure{P|\mathcal P}_3(S)$} des \deleted{8} faces \deleted{\unsure{et}} l'ens
des arêtes \inserted{$A\subset\unsure{P|\mathcal P}_2(S)$} est défini comme l'ens des parties de
card 2 contenues dans une face), par le choix
d'une orientation $\omega$. Choisir cette orientation, revient \inserted{donc}
à la choisir dans le cas standard $X=\mathbb P^1$, $S=\{0,\infty,\pm1,\pm i\}$,
où une face est $\{0,1,i\}$, et les faces \inserted{\uline{orientées}} $\overrightarrow{\{0,1,i\}}$ et
$\overrightarrow{\{0,1,-i\}}$ sont \uline{opposées}. Donc le choix d'une orientation
combinatoire revient au choix d'une racine carrée de
$-1$.

Pour un sous-groupe $V\simeq D_2$ de $G_X$, on désigne
par $\operatorname{Oct}(V)\subset X$ l'octaèdre correspondant, par $I(V)=F/\uline{a}_{\unsure{I|F}}$
l'ens de ses paires de faces antipodiques. Mais : $\Gamma(V)=\operatorname{Nor}_G(V)$
